\section{Limitations and future work}

The characterization in \cref{obj:thm:uniform-frontier} concerns a scalar causal mean at a prespecified latent dose, with fixed public class constants, fixed mean-smoothness and design-decay exponents, a public Gaussian measurement-error scale, and the structural restrictions in \cref{obj:def:model-class}. The proved range \(\sigma\in[0,1/4]\) is an explicit restriction; the analysis does not claim the same uniform constants beyond that cap. The requirement \(\kappa\le2\) is also substantive rather than a normalization: it is used in the dictionary score-rate bounds of \cref{obj:lem:dictionary-score-rate} behind the observable upper bound, and in the comparison with the direct-regression benchmark in \cref{obj:prop:error-free-reduction}. The observable construction and coverage guarantee in \cref{obj:thm:observable-upper,obj:thm:finite-honesty} operate within that specification. These results provide a benchmark for extensions in which model parameters are learned, covariate adjustment is richer, the measurement-error range is widened, or inference covers several doses.

Adaptation over mean smoothness and design decay would require selecting localization weights while accounting for uncertainty about both exponents. The public-score selection in \cref{obj:def:total-estimator} uses the specified class parameters; adaptation would instead require observable criteria that remain valid across candidate classes. Point estimation and honest inference raise distinct questions here. For estimation, the task is to determine the attainable accuracy across classes and any cost of adaptation. For intervals, the task is to establish whether coverage over a union of classes can coexist with shorter expected length on individual classes, and which additional restrictions would permit that improvement.

An unknown measurement-error scale introduces a further source of uncertainty into Gaussian inversion. Replicate measurements or an external calibration sample could supply information about the scale, subject to assumptions linking those measurements to the main sample. An extension would need to establish joint identification, propagate calibration uncertainty through the inverse weights, and compare the required calibration precision with the localization scale. Honest intervals would then have to account for uncertainty from both samples. The transition between measurement-error regimes makes the relative sizes of the calibration and outcome samples a substantive statistical question.

Richer covariates would extend the current two-stratum adjustment to settings with many strata or continuous confounders. Such an analysis would require conditions on covariate support, conditional dose densities, and the regularity of the conditional potential-outcome mean. The statistical problem would combine dose inversion with estimation and aggregation across covariate values. Relevant questions include how covariate dimension affects attainable accuracy, how to control sparsely represented covariate regions, and how uncertainty in covariate adjustment enters an honest interval for the population causal mean.

Simultaneous dose inference would replace evaluation at one prespecified dose with inference over a collection or interval of doses. This extension requires a model describing how weak dose support varies across evaluation points, including its behavior near support boundaries. Uniform approximation and stochastic-error bounds would then be needed for an entire estimated curve. The corresponding questions concern the attainable width of honest confidence bands, the cost of simultaneous coverage, and the relationship between local design decay and uncertainty across doses.

Empirical use also calls for validation of the support and density envelopes. Because the dose is latent, checking these restrictions from contaminated measurements requires an explicit measurement model or additional information. Validation measurements, substantive support restrictions, and sensitivity analysis could help assess plausible specifications. A statistical treatment would need to incorporate uncertainty in estimated support boundaries and envelope constants into estimation and coverage guarantees. Distinguishing assumptions supported by validation data from restrictions supplied by substantive knowledge would make those guarantees interpretable in an application.

The exact public certificates also call for a numerical study separate from the rate theory proved here. A reproducible implementation should evaluate the Fourier and polynomial inverse weights with documented quadrature and stability controls, then report the selected dictionary entry, score, and interval radius over grids of sample sizes and error scales for representative calibrated \(K,\beta,\kappa\). Such diagnostics would measure finite-sample informativeness of the explicit constants without changing the uniform asymptotic and coverage guarantees established in this paper.

Dietary-intake applications are a direction requiring application-specific validation and analysis. A study would need to define the latent intake exposure and causal contrast, assess the measurement-error distribution and its independence conditions, and justify the covariates used for causal adjustment. Replicate assessments or reference measurements could inform the measurement model, while the observed study population would determine the relevant dose-support specification. Connecting these ingredients to the present benchmark requires an analysis of the resulting observed experiment and of uncertainty contributed by exposure calibration and confounder adjustment.

The secondary Gaiffas direct-regression comparison in \cref{obj:prop:error-free-reduction} has a narrower evidentiary status than the paper's own frontier. Its source interface was attested, but an independent source-retrieval check was inconclusive. The comparison is therefore retained as a benchmark for the error-free exponent, with its precise assumptions and formalization scope stated locally at the proposition.
